% Data-driven figures for the final KBS structure.
\begin{figure*}[t]
  \centering
  \includegraphics[width=0.98\textwidth]{figures/fig1_conceptual.pdf}
  \caption{The two conditions studied in this work. A candidate's value changes
  after a similar context is selected. Oracle-Static versus Oracle-Greedy
  measures structural headroom, while cached and response-aware inputs
  determine utility predictability.}
  \label{fig:conceptual}
\end{figure*}

\begin{figure*}[t]
  \centering
  \includegraphics[width=0.98\textwidth]{figures/fig2_headroom.pdf}
  \caption{Redundancy creates state-conditioning headroom. The gap between
  Oracle-Static and Oracle-Greedy grows with duplicate ratio, while MUR
  recovers part of that gap.}
  \label{fig:headroom}
\end{figure*}

\begin{figure*}[t]
  \centering
  \includegraphics[width=0.98\textwidth]{figures/fig3_scaling.pdf}
  \caption{Candidate-set growth creates a ranking bottleneck. Pointwise utility
  MAE remains modest as $K$ increases, while top-1 accuracy falls and
  one-step regret rises.}
  \label{fig:scaling}
\end{figure*}

\begin{figure*}[t]
  \centering
  \includegraphics[width=0.98\textwidth]{figures/fig4_frontier.pdf}
  \caption{Calibrated stopping traces a gain--harm frontier at two harmful
  candidate fractions. The free-threshold control follows a similar frontier.}
  \label{fig:frontier}
\end{figure*}

\begin{figure*}[t]
  \centering
  \includegraphics[width=0.98\textwidth]{figures/fig5_real_observability.pdf}
  \caption{Real-data headroom and utility predictability. Both tasks contain
  oracle state-conditioning headroom. On Traffic, adding expert-response
  summaries improves utility predictability beyond cached context features.}
  \label{fig:real_observability}
\end{figure*}
